% ============================================================================
% Appendix I: Evaluation Environment Construction and Anti-Cherry-Picking
% Justifies the public-benchmark router-stress dataset methodology
% ============================================================================

\section{Evaluation Environment Construction}
\label{app:dataset_construction}

This appendix documents how the public-benchmark evaluation environment is
constructed and why the construction process precludes cherry-picking.
We address six desiderata---accessibility, documentation, explicit environment
definition, bias handling, reproducibility, and fairness to
baselines---motivated by current best practices for datasets and benchmarks
in the applied machine-learning community~\cite{hu2024routerbench,
wang2025barp, li2026llmrouterbench}.

% -------------------------------------------------------------------
\subsection{Public, Accessible Sources}
\label{app:dataset_sources}

All prompts are drawn from publicly available HuggingFace datasets, listed in
Table~\ref{tab:prompt_sources}.
This mirrors how RouterBench~\cite{hu2024routerbench} and
BaRP~\cite{wang2025barp} build evaluation environments from GSM8K, MMLU, ARC,
Natural Questions, and similar public benchmarks.
No private, proprietary, or ad-hoc prompt collections are used; any
researcher with internet access can download the same source data and
reproduce the identical prompt pool.

\begin{table}[h]
\centering
\caption{Public benchmark sources for the evaluation prompt pool.
All datasets are freely available on HuggingFace.}
\label{tab:prompt_sources}
\small
\begin{tabular}{lrl}
\toprule
\textbf{Benchmark} & \textbf{Prompts} & \textbf{Domain} \\
\midrule
MMLU~\cite{hendrycks2021measuring}            & 2{,}672 & Broad academic knowledge \\
GSM8K~\cite{cobbe2021gsm8k}                   & 2{,}328 & Grade-school math \\
HellaSwag~\cite{zellers2019hellaswag}          & 1{,}350 & Commonsense NLI \\
BIG-Bench Hard~\cite{suzgun2022challenging}    & 1{,}541 & Hard reasoning (27 tasks) \\
ARC-Challenge~\cite{clark2018arc}              & 1{,}050 & Science reasoning \\
OpenBookQA~\cite{mihaylov2018suit}             & 1{,}037 & Open-book science QA \\
WinoGrande~\cite{sakaguchi2020winogrande}      & 1{,}005 & Commonsense coreference \\
TruthfulQA~\cite{lin2022truthfulqa}            &    712 & Truthfulness probing \\
MBPP~\cite{austin2021program}                  &    305 & Code generation \\
\midrule
\textbf{Total (after dedup)}                   & \textbf{12{,}000} & \\
\bottomrule
\end{tabular}
\end{table}

% -------------------------------------------------------------------
\subsection{Transparent, Documented Construction Pipeline}
\label{app:dataset_pipeline}

The entire construction pipeline is implemented in a single deterministic
script (\texttt{build\_router\_pareto\_dataset.py}) and proceeds as follows:

\begin{enumerate}[nosep,leftmargin=*]
  \item \textbf{Source loading.}
    Load prompts from each benchmark via the HuggingFace \texttt{datasets}
    library using fixed splits and deterministic iteration order.

  \item \textbf{Quality filtering.}
    Discard prompts with $<$20 or $>$5{,}000 characters, or with
    $<$50\% ASCII content (removes encoding artifacts and non-English prompts).

  \item \textbf{Deduplication.}
    (a)~Exact string dedup against the pool and against any existing
    canonical prompts.
    (b)~Semantic dedup via cosine similarity of
    \texttt{all-MiniLM-L6-v2} embeddings at threshold 0.85, removing
    near-duplicate reformulations across benchmarks.

  \item \textbf{Full-information evaluation.}
    For each prompt, query \emph{all} $K{=}3$ arms (budget, mid-tier,
    premium) via the same OpenRouter API, collect the response, and
    score it with a single DeepSeek-R1 judge using the v3 continuous
    rubric (Reasoning Quality~40\%, Instruction Following~30\%,
    Communication Quality~30\%; \S\ref{sec:data}).
    Record actual input/output token counts and compute exact
    per-request cost from published model pricing.

  \item \textbf{Pareto classification.}
    Classify each prompt based on the reward spread across arms and
    whether a cheaper arm achieves near-best quality
    (Section~\ref{app:pareto_classification}).

  \item \textbf{Output.}
    Write full-information JSONL logs containing prompt text, per-arm
    responses, per-arm rewards, per-arm token counts, and per-arm costs.
\end{enumerate}

Every parameter (random seed, cosine threshold, quality gates, judge
rubric weights) is declared in the script header and can be inspected
or modified by reviewers.

% -------------------------------------------------------------------
\subsection{Explicit Environment Definition, Not Hidden Cherry-Picking}
\label{app:pareto_classification}

A natural concern with any curated evaluation set is whether the
construction process implicitly selects prompts that flatter the
proposed method.
Our construction provides three structural safeguards against this:

\paragraph{Full-information, symmetric evaluation.}
Every arm is evaluated on every prompt \emph{before} any prompt
selection occurs.
The classification rule operates on the full reward--cost matrix,
not on the output of a learned router.
This ensures that any baseline---``always budget,'' ``always premium,''
best single arm, or any supervised router---operates in exactly the
same environment.

\paragraph{Deterministic, inspectable classification.}
Prompts are labelled by a simple rule applied uniformly:
\begin{itemize}[nosep,leftmargin=*]
  \item \textbf{Trivial:} reward spread $\leq 0.05$ (all arms perform
    similarly; routing adds no value).
  \item \textbf{Pareto-interesting:} reward spread $\geq 0.10$ \emph{and}
    a cheaper arm achieves within 0.05 of the best reward (a genuine
    cost--quality tradeoff exists).
  \item \textbf{Hard-but-dominated:} high spread, but the most expensive
    arm strictly dominates (routing cannot save cost without quality loss).
  \item \textbf{Moderate:} intermediate cases.
\end{itemize}
These thresholds are published in the code and can be freely adjusted
by reviewers to test sensitivity.
RouterBench~\cite{hu2024routerbench} similarly defines its evaluation
in the cost--quality plane and constructs AIQ scores from the full
Pareto frontier; our classification is a lightweight analogue that makes
the ``routing difficulty'' of each prompt explicit rather than
implicit.

\paragraph{Balanced sampling, not selective sampling.}
The final evaluation set retains \emph{all} Pareto-interesting prompts
and supplements them with proportional samples from trivial and
hard-but-dominated bins.
This produces a ``router stress-test'' environment that over-represents
prompts where routing matters, while still including prompts where it
does not---ensuring that aggregate metrics are not inflated by
excluding easy-for-the-router cases.

% -------------------------------------------------------------------
\subsection{Bias and Difficulty Are Handled Explicitly}
\label{app:difficulty_bias}

Existing router evaluations typically inherit whatever difficulty
distribution the upstream benchmarks impose.
BaRP~\cite{wang2025barp}, for example, splits each task 80/20 and mixes
them; RouterBench~\cite{hu2024routerbench} aggregates across tasks
without stratifying by router-relevant difficulty.
In both cases, the implicit difficulty distribution is uncontrolled and
unreported.

Our construction instead:
\begin{enumerate}[nosep,leftmargin=*]
  \item Uses the full-information log to compute how much arms disagree
    per prompt---a data-driven measure of routing-relevant difficulty.
  \item Explicitly defines difficulty strata and states the sampling
    distribution over those strata.
  \item Provides \emph{both} the natural (unstratified) and
    difficulty-balanced environments so that results can be reported
    on each.
\end{enumerate}

This is not an attempt to ``game'' the difficulty mix; it is an explicit
declaration that we are constructing a stress-test environment and
documenting how.
The natural environment is always available for comparison, and the
stratification rule is public.

% -------------------------------------------------------------------
\subsection{Reproducible Evaluation in the RouterBench Style}
\label{app:reproducibility}

RouterBench formalises router evaluation as comparing points in the
cost--quality plane and constructing Pareto frontiers and AIQ
scores~\cite{hu2024routerbench}.
Our script produces a JSONL log with:
\begin{itemize}[nosep,leftmargin=*]
  \item The prompt text and benchmark source.
  \item Per-arm responses, actual token counts, and computed costs.
  \item Per-arm reward scores from the R1 judge (with full rubric
    sub-scores).
  \item A difficulty label enabling ``router-stress'' vs.\ ``natural''
    evaluation splits.
\end{itemize}
With this file and the construction script, any group can: (a)~recompute
the same cost--quality curves; (b)~change classification thresholds and
regenerate the dataset; (c)~substitute their own judge or reward
function and re-evaluate; or (d)~add arms and extend the environment.
This level of artifact transparency satisfies modern reproducibility
standards for both research papers and datasets/benchmarks
submissions.

% -------------------------------------------------------------------
\subsection{Fairness to Baselines}
\label{app:baseline_fairness}

Crucially, the prompt classification is performed \emph{without}
reference to any learned routing policy.
Prompts are labelled solely by the arm-level reward--cost matrix,
which is computed symmetrically for all arms.
This means:

\begin{itemize}[nosep,leftmargin=*]
  \item Static baselines (``always budget,'' ``always premium,''
    best-on-average arm) naturally inhabit the same environment
    and can be plotted on the same cost--quality Pareto frontier.
  \item Supervised baselines (KNN, SVM, MLP routers) operate on
    identical features and reward signals.
  \item The oracle upper bound is computed from the same full-information
    log.
\end{itemize}

No information leaks from the bandit's learned policy into the
environment construction.
If the router underperforms on the stress-test slice, the environment
will reveal this just as readily as it reveals gains---unlike a
cherry-picked set, which would hide failure modes by construction.

% -------------------------------------------------------------------
\subsection{Reward Transparency}
\label{app:reward_transparency}

The reward signal for each (prompt, arm) pair is produced by a
DeepSeek-R1 judge~\cite{deepseekai2025deepseekr1} using a structured
chain-of-thought rubric with three continuous dimensions:
\emph{Reasoning Quality}~(weight 0.40), \emph{Instruction
Following}~(0.30), and \emph{Communication Quality}~(0.30).
Each dimension is scored on a continuous $[0, 1]$ scale with
calibrated anchor descriptions at five levels (0.0--0.2 through
0.9--1.0), and the final reward is the weighted sum.

The rubric text, parsing logic, and score aggregation are implemented
in the same codebase and can be audited.
This addresses the concern that reward definitions in router
evaluations are often opaque: the judge model, prompt template, and
scoring weights are all public and deterministic (\texttt{temperature}=0).
